\subsection{No supercritical V role}

\begin{proposition}[One-gap common-disk closure]
\label{prop:readable-nplus-zero-one-gap}
Assume $N_{\mathrm{gap}}=1$, $N_+=0$, no V role is Vd1 or Vd2, and at most
two V roles are T3-like.  Then the cover is impossible.
\end{proposition}

\begin{proof}
Normalize the gap to $J=[X(\ell),X(r)]$ and put
$p=1-r$, $q=\ell$, $c_*=c_{\max}(p,q)$.  Lemma
\ref{lem:readable-one-gap-common-pair} supplies the common pair.  Type-aware
radial forcing and common-pair domination put all six points
$(1-c_*)V_i$ in $U_C$, including the permitted T3-like neighboring terms.
Their convex hull contains $\mathcal D_{h(1-c_*)}$, while the actual gap lies
in $U_C$.  Theorem~\ref{thm:new-complementary-gap} contradicts
Lemma~\ref{lem:new-compact-open-shrink}.
\end{proof}

\CaseCard{NG0}{${N_{\rm gap}=1}$, $N_+=0$, no Vd1/Vd2, $t\le2$.}
{$\mathcal D_{h(1-c_*)}\cup J(p,q)$.}
{Theorem~\ref{thm:new-complementary-gap}.}
{Routing rows R7 and R9 with one gap.}

\begin{proposition}[Two-gap common-pair closure]
\label{prop:readable-two-gap-common-closure}
Assume CE2, both center traces contain V-gaps, and the intervening roles
$T_1,\ldots,T_5$ are nonsupercritical Vd0 or T3-like roles.  Then the cover
is impossible, independently of the criticality of $T_0$.
\end{proposition}

\begin{proof}
Lemma~\ref{lem:readable-two-gap-common-pair} supplies
$(p,q)=(W-\alpha,R-\delta)$.  Common-pair domination controls every
permitted T3-like neighboring term, so the type-aware witnesses $D_2,D_4$
belong to $U_C$.  The shortened proof of
Theorem~\ref{thm:new-ce2-short-ray} places one strictly beyond the
corresponding C exit.
\end{proof}

\CaseCard{NG1}{CE2, two gaps, $N_+\in\{0,1\}$, intervening roles Vd0 or
T3-like.}{$D_2,D_4$ and the two actual gaps.}
{Theorem~\ref{thm:new-ce2-short-ray}.}
{The two-gap portions of routing rows R7, R9, R10, and R11.}

\begin{corollary}[All-Vd0 and T3-like $N_+=0$ gap closures]
\label{prop:new-nplus-zero-gap-closures}
Every nonzero-gap state with $N_+=0$, no Vd1/Vd2 role, and at most two
T3-like roles is impossible.
\end{corollary}

\begin{proof}
There is one or two gaps.  Apply
Proposition~\ref{prop:readable-nplus-zero-one-gap} or
Proposition~\ref{prop:readable-two-gap-common-closure}.
\end{proof}

\subsection{One gap, one supercritical role, and all V roles Vd0}

Let the gap be $[X(\ell),X(r)]$.  Let $C_i$ be the actual own-radial reach
of $T_i$ and put
\[
 P_i=(1-C_i)V_i.
\]

\begin{lemma}[The transverse set is center-forced]
\label{lem:readable-k410-forced}
Every $P_i$ is missed by all six open V roles and lies in $U_C$.  In
particular,
\[
 K_{\rm tr}
 =
 \{O,M_0,X(\ell),X(r),P_2,P_3,P_4\}
 \subset U_C.
 \tag{6.35}
\]
\end{lemma}

\begin{figure}[!htbp]
\centering
\begin{minipage}[t]{.325\linewidth}
\centering
\resizebox{\linewidth}{!}{%
\begin{tikzpicture}[scale=2.05,>=Latex]
  \node[panellabel,anchor=west] at (-1.08,1.18) {(a) the finite set $K_{410}$};
  \HexCoords
  \DrawHexagon
  \DrawRays
  \draw[gap] ($(V0)!.20!(V1)$)--($(V0)!.34!(V1)$);
  \def\rr{{.13,.19,.15,.10,.21,.16}}
  \foreach \i in {0,...,5}{
    \pgfmathparse{\rr[\i]}
    \coordinate (P\i) at ($(O)!{\pgfmathresult}!(V\i)$);
    \draw[VertexOrange,fill=white,line width=.65pt] (P\i) circle (.027);
    \fill[DemandRed] (P\i) circle (.012);
  }
  \fill[CenterBlue] (O) rectangle +(0.035,0.035);
  \fill[CenterBlue] ($(O)!.5!(V0)$) rectangle +(0.035,0.035);
  \fill[GapPurple] ($(V0)!.20!(V1)$) -- ++(.03,.02) -- ++(-.03,.02) -- ++(-.03,-.02) -- cycle;
  \fill[GapPurple] ($(V0)!.34!(V1)$) -- ++(.03,.02) -- ++(-.03,.02) -- ++(-.03,-.02) -- cycle;
  \node[figlabel,below left] at (O) {$O$};
  \node[figlabel,above] at ($(O)!.5!(V0)$) {$M_0$};
  \node[figlabel,right] at (P0) {$P_0=(1-C_0)V_0$};
  \node[figlabel,align=center] at (0,-1.18)
    {circles: actual radial endpoints $P_i$\\diamonds: $X(\ell),X(r)$};
\end{tikzpicture}%
}
\end{minipage}\hfill
\begin{minipage}[t]{.35\linewidth}
\centering
\resizebox{\linewidth}{!}{%
\begin{tikzpicture}[scale=2.05,>=Latex]
  \node[panellabel,anchor=west] at (-1.08,1.18) {(b) a candidate center triangle};
  \HexCoords
  \DrawHexagon
  \DrawRays
  \draw[centerrole]
    (-.58,-.32)--(.88,-.20)--(-.02,.86)--cycle;
  \coordinate (P0) at ($(O)!.13!(V0)$);
  \coordinate (E0) at ($(O)!.30!(V0)$);
  \draw[CenterBlue,line width=1pt] (O)--(E0);
  \draw[VertexOrange,line width=1.2pt] (V0)--(P0);
  \draw[VertexOrange,fill=white,line width=.65pt] (P0) circle (.027);
  \fill[DemandRed] (P0) circle (.012);
  \fill[CenterBlue] (E0) circle (.023);
  \node[figlabel,above] at (P0) {$P_i$};
  \node[figlabel,below] at (E0) {$E_i'=d_i'V_i$};
  \draw[dimline] ($(O)+(0,-.13)$)--($(P0)+(0,-.13)$)
    node[midway,below=1pt,figlabel] {$1-C_i$};
  \draw[dimline] ($(O)+(0,.13)$)--($(E0)+(0,.13)$)
    node[midway,above=1pt,figlabel] {$d_i'$};
  \node[figlabel,align=center] at (0,-1.18)
    {$P_i\in T_C'\Longrightarrow d_i'\ge1-C_i$\\
     equivalently $C_i\ge1-d_i'$};
\end{tikzpicture}%
}
\end{minipage}\hfill
\begin{minipage}[t]{.285\linewidth}
\centering
\resizebox{\linewidth}{!}{%
\begin{tikzpicture}[x=5cm,y=1cm,>=Latex]
  \node[panellabel,anchor=west] at (-.03,2.62) {(c) closing a boundary gap};
  \node[figlabel,anchor=west] at (0,2.14) {before};
  \draw[line width=.65pt] (0,1.70)--(1,1.70);
  \draw[actualreach] (0,1.76)--(.38,1.76);
  \draw[centertrace] (.53,1.64)--(.77,1.64);
  \draw[gap] (.38,1.70)--(.53,1.70);
  \node[figlabel,text=GapPurple] at (.455,1.98) {gap};
  \node[figlabel,anchor=west] at (0,1.08) {after $C_i\ge1-d_i'$};
  \draw[line width=.65pt] (0,.64)--(1,.64);
  \draw[actualreach] (0,.70)--(.58,.70);
  \draw[centertrace] (.53,.58)--(.77,.58);
  \draw[RadialTeal,line width=1pt] (.53,.64)--(.58,.64);
  \node[figlabel,text=RadialTeal] at (.555,.91) {overlap};
  \node[figlabel,align=center] at (.5,-.03)
    {every candidate gap closes\\$\Rightarrow$ CE1 or CE2 terminal};
\end{tikzpicture}%
}
\end{minipage}
\caption{The actual-reach mechanism behind the anisotropic set $K_{410}$.
The six points $P_i=(1-C_i)V_i$ retain the actual own-radial reaches rather
than a common relaxed radius.  If a candidate center triangle $T_C'$ contains
$K_{410}$, its exit $d_i'$ on every radial arm satisfies
$d_i'\ge1-C_i$, equivalently $C_i\ge1-d_i'$.  The corresponding V trace then
meets the center contribution and closes the associated boundary gap, forcing
the candidate into the CE1 or CE2 terminal used in the proof.  Lengths are
schematic; the order relations are exact.}
\label{fig:k410-actual-reach-certificate}
\end{figure}


\begin{proof}
The point $P_i$ is the closed own-trace endpoint and is not in $U_i$ by
Lemma~\ref{lem:new-open-trace-endpoint}.  Adjacent roles are Vd0 and the
nonlocal roles are excluded by diameter.  The actual gap and all $P_i$ are
therefore center-forced.
\end{proof}

\begin{lemma}[Endpoint upper squeeze]
\label{lem:readable-k410-upper-bound}
Give an open unit triangle $U'$ containing $K_{\rm tr}$ the signed center
variables, and put
\[
 k=\eta+\alpha+\delta,\qquad
 X=R-\delta,\qquad Q=\frac{k}{2R}.
\]
Then the unique supercritical role satisfies
\[
 \boxed{A_0<1-Q.}
 \tag{6.36}
\]
\end{lemma}

\begin{proof}
Since the two actual gap endpoints lie in the open trace of $U'$,
\[
 B_0=\ell>\frac{k}{R}=2Q,
 \qquad
 A_1=1-r>X.
\]
The two boundary endpoints of $T_0$ have distance at most one:
\[
 A_0^2+A_0B_0+B_0^2\le1.
\]
Hence, with the decreasing zero-radial diameter envelope,
\[
 A_0
 \le M_0(B_0)
 <M_0(2Q)
 =-Q+\sqrt{1-3Q^2}
 <1-Q.
\]
The old radial point $P_0$ is not used.
\end{proof}

\begin{proposition}[CE1 three-transverse return certificate]
\label{prop:new-ce1-direct-certificate}
Use the signed CE1 center variables
\[
0<R<1,
\qquad w=1-R,
\qquad E=\sqrt{1-Rw},
\]
\[
\eta=1-E,
\qquad P=E(1-E),
\]
and write
\[
 A=\alpha,
 \qquad D=\delta.
\]
Assume
\[
 A>0,
 \qquad D>0,
 \qquad A+wD<P,
 \qquad D+RA\ge P.
 \tag{6.12}
\]
Put
\[
 X=R-D,
 \qquad m=\frac AR,
 \qquad h_0=\frac{\eta+A+D}{2R}.
 \tag{6.13}
\]
Let $T_0$ be the unique supercritical role and let
$T_1,\ldots,T_5$ be nonsupercritical $\Vdzero$ roles.  Assume
\[
 A_1\ge X,
\]
\[
 B_i+A_{i+1}\ge1\quad(1\le i\le4),
 \qquad B_5+A_0\ge1,
\]
and only the three transverse radial reaches
\[
 C_4\ge1-A,\qquad C_3\ge1-m,\qquad C_2\ge1-D.
 \tag{6.14}
\]
Then
\[
 \boxed{A_0>1-h_0.}
 \tag{6.15}
\]
\end{proposition}

\begin{proof}
The direct path $T_5\to T_4\to T_3\to T_2\to T_1$, its two selected-chord
steps, and the final low-root contradiction are proved in
Appendix~\ref{subsec:calc-ce1-reverse-path}.  That calculation uses precisely
the five displayed boundary handoffs and the radial reaches at
$T_4,T_3,T_2$; no reach on $r_0,r_1,r_5$ and no composed map is used.
\end{proof}


\begin{proposition}[CE2 transverse threshold terminal]
\label{prop:readable-k410-ce2}
No CE2 open unit triangle contains $K_{\rm tr}$.
\end{proposition}

\begin{proof}
Let $U'$ be such a triangle.  Its gap trace gives $A_1>X$, while
Lemma~\ref{lem:readable-k410-upper-bound} gives $A_0<1-Q$.
The five nonsupercritical roles and the V-gap-free path imply
\[
 A_1<A_2<A_3<A_4<A_5<A_0.
\]
Consequently
\[
 B_2>Q,\qquad B_4>Q.
 \tag{6.37}
\]
Containment of $P_2,P_4$ gives
\[
 C_2>1-\delta,\qquad C_4>1-\alpha.
 \tag{6.38}
\]
Appendix~\ref{subsec:calc-ce2-one-gap-threshold} proves
\[
 X>\epsilon(\alpha),
 \qquad
 \min\{\epsilon(\alpha),\epsilon(\delta)\}<Q.
 \tag{6.39}
\]
If $\epsilon(\alpha)<Q$, the direct threshold at $T_4$ gives
$B_4\le\epsilon(\alpha)<Q$, contradicting (6.37).  Otherwise
$\epsilon(\delta)<Q$ and
\[
 A_2>X>\epsilon(\alpha)\ge Q>\epsilon(\delta);
\]
the threshold at $T_2$ gives $B_2\le\epsilon(\delta)<Q$.  Thus no CE2
candidate exists.  The point $P_3$ is not needed in this branch.
\end{proof}

\begin{proposition}[CE1 three-transverse return]
\label{prop:readable-k410-ce1}
No CE1 open unit triangle contains $K_{\rm tr}$.
\end{proposition}

\begin{figure}[!htbp]
\centering
\resizebox{.92\linewidth}{!}{%
\begin{tikzpicture}[
  x=1cm,y=1cm,>=Latex,
  state/.style={draw=black!38,rounded corners=1.5pt,fill=black!2,
    minimum height=.78cm,text width=2.35cm,align=center,font=\scriptsize},
  hard/.style={draw=RadialTeal,rounded corners=1.5pt,fill=RadialTeal!7,
    minimum height=.78cm,text width=2.35cm,align=center,font=\scriptsize},
  arr/.style={-{Stealth[length=1.8mm]},draw=black!60,line width=.65pt},
  harr/.style={-{Stealth[length=1.8mm]},draw=RadialTeal,line width=.85pt}
]
  \node[panellabel,anchor=west] at (-.1,2.55)
    {CE1 reverse path: from a deficient terminal to the final contradiction};

  \node[state] (a0) at (0,1.35) {$A_0\le1-h_0$};
  \node[state] (b5) at (3.05,1.35) {$B_5\ge h_0$};
  \node[state] (b4) at (6.10,1.35) {$B_4\ge h_0$};
  \node[hard]  (b3) at (9.15,1.35) {$B_3>L_1$};

  \node[hard]  (b2) at (0,-.55) {$B_2>L_2>\epsilon(D)$};
  \node[hard]  (a2) at (3.05,-.55) {$A_2\le\epsilon(D)$};
  \node[state] (b1) at (6.10,-.55) {$B_1>1-X$};
  \node[state,draw=DemandRed,fill=DemandRed!7]
    (contra) at (9.15,-.55)
    {$A_1\ge X$, $A_1+B_1\le1$\\contradiction};

  \draw[arr] (a0)--node[above=3pt,figlabel,fill=white,inner sep=.6pt] {companion edge} (b5);
  \draw[arr] (b5)--node[above=3pt,figlabel,fill=white,inner sep=.6pt] {Vd0 monotonicity} (b4);
  \draw[harr] (b4)--node[above=3pt,figlabel,fill=white,inner sep=.6pt] {selected $Q_+$ at $T_4$} (b3);
  \draw[harr] (b3.east)--++(.55,0)--++(0,-1.10)--
    node[right=2pt,figlabel,align=left,fill=white,inner sep=.6pt] {selected $Q_+$ at $T_3$\\and the scalar estimate}
    ++(-10.25,0)--(b2.west);
  \draw[harr] (b2)--node[above=3pt,figlabel,fill=white,inner sep=.6pt] {threshold at $T_2$} (a2);
  \draw[arr] (a2)--node[above=3pt,figlabel,fill=white,inner sep=.6pt] {coverage of $e_{1,2}$} (b1);
  \draw[arr,draw=DemandRed] (b1)--(contra);

  \node[figlabel,atlasbox,align=left,anchor=west] at (1.15,-1.68)
    {grey arrows: edge coverage or nonsupercriticality;\qquad
     teal arrows: the genuinely local high-radial steps};
  \node[figlabel,align=center] at (5.25,-2.33)
    {$L_1$ and $L_2$ are direct selected-arc lower bounds; no composed reach map is introduced.};
\end{tikzpicture}%
}
\caption{Logical structure of the CE1 direct reverse-path certificate.  The
proof starts from the negation of the desired terminal bound at $T_0$, moves
through the actual boundary reaches, uses two selected-arc estimates and one
high-radial threshold, and returns to $T_1$ with incompatible incoming and
outgoing demands.}
\label{fig:ce1-reverse-path}
\end{figure}


\begin{proof}
For a CE1 candidate, containment of $P_4,P_3,P_2$ gives the three radial
demands
\[
 C_4>1-\alpha,\qquad
 C_3>1-\frac{\alpha}{R},\qquad
 C_2>1-\delta.
\]
These are exactly the only radial hypotheses used by
Proposition~\ref{prop:new-ce1-direct-certificate}.  Its reverse path through
$T_5,T_4,T_3,T_2,T_1$ proves
\[
 A_0>1-Q,
\]
whereas Lemma~\ref{lem:readable-k410-upper-bound} gives $A_0<1-Q$.
The former points $P_0,P_1,P_5$ do not occur in the certificate.
\end{proof}

\begin{proposition}[Transverse seven-point enclosure]
\label{prop:new-nplus-one-all-vd0}
The compact set
\[
 K_{\rm tr}
 =
 \{O,M_0,X(\ell),X(r),P_2,P_3,P_4\}
\]
has
\[
 \boxed{\Lambda(K_{\rm tr})\ge1.}
 \tag{6.40}
\]
Consequently the former superset
\[
 K_{410}=\{O,M_0,X(\ell),X(r),P_0,\ldots,P_5\}
\]
also has $\Lambda(K_{410})\ge1$.
\end{proposition}

\begin{proof}
An open unit triangle containing $K_{\rm tr}$ has a positive boundary trace
and is CE1 or CE2.  Propositions~\ref{prop:readable-k410-ce1} and
\ref{prop:readable-k410-ce2} exclude both alternatives.  Apply
Lemma~\ref{lem:new-compact-open-shrink} and monotonicity of $\Lambda$.
\end{proof}

\CaseCard{NG2}{One gap, $N_+=1$, all V roles Vd0.}
{$K_{\rm tr}=\{O,M_0,X(\ell),X(r),P_2,P_3,P_4\}$.}
{Endpoint upper squeeze versus CE1/CE2 transverse return.}
{Routing row R10 with one gap.}

\subsection{A common neighboring-midpoint rescuer budget}

\begin{lemma}[Rescuer endpoint and boundary-tail budget]
\label{lem:readable-rescuer-tail-budget}
Suppose a special role at $V_0$ has boundary endpoint $a$ on $e_{5,0}$ and
supported interval $[c,u]\subset r_1$, measured from $V_1$ toward $O$.
Put
\[
 \varepsilon=1-u,\qquad M=M_c^{\rm sup}.
\]
Assume
\[
 \varepsilon V_1\in U_C,\qquad a+\varepsilon\le1,
\]
\[
 a\le1-M,\qquad
 \frac{a}{a+\varepsilon}\le1-M.
 \tag{6.41}
\]
If $T_1$ is uniquely supercritical with own-radial reach at least $c$ and
$T_2,T_3,T_4,T_5$ are nonsupercritical on the center-free four-edge path,
then the skeleton is not covered.
\end{lemma}

\begin{proof}
Let $h_T$ be the far boundary demand on $T_5$.  If the companion center trace
does not hide $a$, then $h_T\ge1-a\ge M$.

In the only hiding configuration, its near endpoint gives $k/W\le a$, while
$\varepsilon V_1\in U_C$ gives $\delta/R\ge\varepsilon$.  Hence
\[
 Wa\ge k=\eta+\alpha+\delta>\alpha+R\varepsilon,
\]
so
\[
 a-R(a+\varepsilon)>\alpha.
\]
Also $Wa>\delta\ge R\varepsilon$, so
$R<a/(a+\varepsilon)$.  Since $a+\varepsilon\le1$,
\[
 R+\alpha
 <a+R(1-a-\varepsilon)
 <\frac{a}{a+\varepsilon}
 \le1-M.
\]
Thus again $h_T\ge M$.

The strict-supercritical envelope gives $B_1<M$.  Adding
\[
 A_2\ge1-B_1,\quad
 B_2+A_3\ge1,\quad B_3+A_4\ge1,\quad B_4+A_5\ge1,
 \quad B_5\ge h_T
\]
yields
\[
 \sum_{i=2}^5(A_i+B_i)\ge4+h_T-B_1>4,
\]
contrary to nonsupercriticality.
\end{proof}

\subsection{Exactly one T3-like role}

\begin{lemma}[The T3-like O-side endpoint is center-forced]
\label{lem:readable-t3-o-side-endpoint}
Normalize the unique T3-like role to $T_0$, its supported arm to $r_1$, and
the unique supercritical role to $T_1$.  If the supported interval is
$[c,u]$ with $c\le1/2\le u$, then
\[
 P_T=(1-u)V_1\in U_C.
\]
\end{lemma}

\begin{figure}[!htbp]
\centering
\begin{minipage}[t]{.47\linewidth}
\centering
\resizebox{.96\linewidth}{!}{%
\begin{tikzpicture}[x=6cm,y=1cm,>=Latex]
  \node[panellabel,anchor=west] at (-.02,2.30) {(a) the supported interval on $r_1$};
  \draw[flow] (0,1.15)--(1,1.15);
  \fill (0,1.15) circle (.017) node[left=2pt,figlabel] {$V_1$};
  \fill (1,1.15) circle (.017) node[right=2pt,figlabel] {$O$};
  \coordinate (C) at (.24,1.15);
  \coordinate (M) at (.50,1.15);
  \coordinate (U) at (.73,1.15);
  \draw[actualreach] (C)--(U);
  \draw[VertexOrange,fill=white,line width=.7pt] (U) circle (.025);
  \fill[DemandRed] (U) circle (.011);
  \fill[black] (C) circle (.019) node[above=3pt,figlabel] {$c$};
  \fill[CenterBlue] (M) rectangle +(0.026,0.026);
  \node[above=3pt,figlabel] at (M) {$M_1$};
  \node[above=3pt,figlabel] at (U) {$u$};
  \draw[dimline] ($(U)+(0,-.22)$)--($(1,1.15)+(0,-.22)$)
    node[midway,below=1pt,figlabel] {$1-u$};
  \node[figlabel,align=center] at (.52,.12)
    {$P_T=(1-u)V_1$ is the O-side endpoint\\
     coordinate on the top axis is measured from $V_1$ toward $O$};
\end{tikzpicture}%
}
\end{minipage}\hfill
\begin{minipage}[t]{.47\linewidth}
\centering
\resizebox{.96\linewidth}{!}{%
\begin{tikzpicture}[scale=2.15,>=Latex]
  \node[panellabel,anchor=west] at (-1.08,1.22) {(b) why $P_T$ is center-forced};
  \HexCoords
  \DrawHexagon
  \DrawRays
  \coordinate (PT) at ($(O)!.27!(V1)$);
  \coordinate (M1) at ($(O)!.5!(V1)$);
  \draw[vertexrole]
    (V0)--($(V0)!.76!(V1)$)--($(O)!.20!(V1)$)--cycle;
  \draw[DemandRed,dashed,line width=.9pt]
    (V1)--($(V1)!.48!(V0)$)--($(O)!.48!(V1)$)--cycle;
  \draw[actualreach] ($(V1)!.22!(O)$)--($(V1)!.73!(O)$);
  \fill[DemandRed] (PT) circle (.025) node[left=2pt,figlabel] {$P_T$};
  \fill[CenterBlue] (M1) rectangle +(0.025,0.025);
  \node[figlabel,right] at (M1) {$M_1$};
  \draw[flow] (PT)--($(PT)!1.55!(O)$);
  \node[figlabel,align=left,anchor=west] at (-.96,-.72)
    {$T_0$: open endpoint at $P_T$\\
     $T_1$: misses $M_1$, hence the O-side segment\\
     other V roles: excluded by type or diameter};
  \node[figlabel,atlasbox,align=center] at (.48,-1.05)
    {$P_T\in U_C$};
\end{tikzpicture}%
}
\end{minipage}
\caption{Construction of the T3-like rescuer endpoint.  The interval
$[c,u]$ is parametrized from $V_1$ toward $O$, so its O-side endpoint is
$P_T=(1-u)V_1$.  The T3-like open role omits its endpoint, the adjacent
supercritical role contains $V_1$ but misses $M_1$ and therefore misses the
O-side segment by convexity, and all remaining V roles are excluded.  Thus
$P_T$ is forced into the C triangle.}
\label{fig:t3-rescuer-endpoint}
\end{figure}


\begin{proof}
The O-side endpoint is not in the open T3-like role.  The adjacent
supercritical role contains $V_1$ but misses $M_1$, and the remaining roles
are Vd0 or nonlocal.
\end{proof}

\begin{proposition}[One-gap T3-like rescuer terminal]
\label{prop:readable-one-t3-one-gap}
Assume $N_+=1$, exactly one V role is T3-like, no role is Vd1/Vd2, and
there is exactly one V-gap.  Then the skeleton is not covered.
\end{proposition}

\begin{proof}
Use Lemma~\ref{lem:readable-t3-o-side-endpoint}.  Let $a$ be the T3-like
boundary reach on $e_{5,0}$.  In the translated normal form put
\[
 R_0=\frac{D+\sqrt{4-3D^2}}2,\qquad
 x=\frac{aD}{R_0},\qquad
 \theta=\frac{D-1}{R_0}.
\]
The midpoint constraints give
\[
 \frac{1-4\theta+\theta^2}{2(1-2\theta)}
 \le x\le\frac12-\theta.
\]
The quadratic calculation in Appendix~\ref{subsec:calc-t3-rescuer} proves,
with $\varepsilon=1-u$ and $M=M_c^{\rm sup}$,
\[
 a\le1-M,\qquad
 \frac{a}{a+\varepsilon}\le1-M.
\]
Moreover $a+\varepsilon=R_0/D\le1$.  Apply
Lemma~\ref{lem:readable-rescuer-tail-budget}.
\end{proof}

\begin{proposition}[Two-gap T3-like terminal]
\label{prop:readable-one-t3-two-gap}
Under the same vertex-type hypotheses, a two-gap state is impossible.
\end{proposition}

\begin{proof}
Apply Proposition~\ref{prop:readable-two-gap-common-closure}; common-pair
domination bounds the permitted T3-like neighboring capacity.
\end{proof}

\begin{corollary}[Direct T3-like rescuer witness]
\label{prop:new-one-t3-terminal}
Assume $N_+=1$, exactly one V triangle is T3-like, no V triangle is
Vd1/Vd2, and at least one V-gap is present.  Then the skeleton is not covered.
\end{corollary}

\begin{proof}
Use Proposition~\ref{prop:readable-one-t3-one-gap} or
Proposition~\ref{prop:readable-one-t3-two-gap}.
\end{proof}

\CaseCard{NG3}{$N_+=1$, exactly one T3-like role, no Vd1/Vd2.}
{$P_T$ and the common rescuer budget in one gap; $D_2,D_4$ in two gaps.}
{Lemma~\ref{lem:readable-rescuer-tail-budget} or the CE2 short-ray theorem.}
{Routing row R11.}

\subsection{Exactly one Vd1/Vd2 role in CE2}

Throughout let $T_\sigma$ be the unique supercritical role and $T_\tau$ the
unique Vd1/Vd2 role.  If a further role has positive adjacent support,
Method~1 applies; otherwise every other role is Vd0.

\begin{proposition}[Adjacent one-Vd radial separation]
\label{prop:readable-one-vd-adjacent}
Assume $\sigma=0$ and, after reflection, $\tau=1$.  Then a point of $r_2$
is missed by all seven roles.
\end{proposition}

\begin{proof}
Let $\rho_R,\rho_L$ be the two residual boundary reaches.  The exact residual
calculation gives
\[
 \rho_R+\rho_L<\frac12,\qquad 4\delta<\rho_L.
\]
The Vd supported endpoint satisfies
\[
 u_{1\to2}<1-\rho_L<1-\delta.
\]
At the ordinary role $T_2$,
\[
 A_2>\frac12+\rho_R,\qquad B_2\ge\rho_L.
\]
Lemma~\ref{lem:new-one-third-radial-envelope}, with
$M=1/2+\rho_R$ and $m=\rho_L$, gives
\[
 C_2<1-\frac{\rho_L}{3}<1-\delta.
\]
Neither local V trace reaches the center interval on $r_2$; all other roles
are excluded by Vd0 locality and diameter.
\end{proof}

\begin{proposition}[Nonadjacent one-Vd radial separation]
\label{prop:readable-one-vd-nonadjacent}
Assume $\sigma=0$ and $\tau\in\{2,3,4\}$.  Then
\[
 D_\tau=\min\{\rho_R,\rho_L\}V_\tau
\]
is missed by all seven roles.
\end{proposition}

\begin{proof}
Appendix~\ref{subsec:calc-nonadjacent-vd} proves
\[
 \alpha+\delta<\min\{\rho_R,\rho_L\}.
\]
Every relevant C exit is at most $\alpha+\delta$, while
Lemma~\ref{lem:readable-vd-corner-margins} puts the Vd own trace strictly
short of $D_\tau$.  Adjacent roles are Vd0 and nonlocal roles fail by
diameter.
\end{proof}

\begin{proposition}[Vd1 neighboring-midpoint rescuer]
\label{prop:readable-vd1-midpoint-rescuer}
Assume $\tau=0$, $T_0$ is Vd1, and after reflection $\sigma=1$.  Then the
four ordinary path roles have total boundary demand greater than four.
\end{proposition}

\begin{proof}
Write the supported interval as $[c,u]$, put
$P=(1-u)V_1$, $\varepsilon=1-u$, and $M=M_c^{\rm sup}$.  The endpoint is
center-forced.  Appendix~\ref{subsec:calc-vd1-rescuer} proves
\[
 a\le1-M,\qquad
 \frac{a}{a+\varepsilon}\le1-M.
\]
The Vd half-unit cap and $u\ge1/2$ give $a+\varepsilon<1$.  Apply
Lemma~\ref{lem:readable-rescuer-tail-budget}.
\end{proof}

\begin{proposition}[Two-chart replacement and output-gap rerouting]
\label{prop:readable-one-vd-replacement}
Assume $\sigma\ne0$, $\tau\ne0$, and the exceptional role is Vd1.  The
corrected two-chart replacement produces six nonsupercritical Vd0 roles
preserving the covered skeleton.  Every recomputed output gap rank is
impossible.
\end{proposition}

\begin{proof}
Midpoint rescue makes the distinguished roles adjacent.  The replacement
uses their two correct vertex charts and makes no claim about preservation of
the input gap rank.  Rank zero is excluded by the boundary-complete length
theorem; rank one by the common-disk theorem; rank two by the shortened CE2
short-ray theorem.
\end{proof}

\begin{figure}[!htbp]
\centering
\begin{minipage}[t]{.485\linewidth}
\centering
\resizebox{.96\linewidth}{!}{%
\begin{tikzpicture}[x=6cm,y=1cm,>=Latex]
  \node[panellabel,anchor=west] at (-.02,2.28)
    {(a) adjacent Vd role: an uncovered interval on $r_2$};
  \draw[flow] (0,1.12)--(1,1.12);
  \fill (0,1.12) circle (.017) node[left=2pt,figlabel] {$V_2$};
  \fill (1,1.12) circle (.017) node[right=2pt,figlabel] {$O$};
  \coordinate (U) at (.46,1.12);
  \coordinate (C) at (.58,1.12);
  \coordinate (Q) at (.74,1.12);
  \draw[VertexOrange,line width=1.1pt] (0,1.19)--(U);
  \draw[RadialTeal,line width=1.1pt] (0,1.05)--(C);
  \draw[centertrace] (Q)--(1,1.12);
  \fill[VertexOrange] (U) circle (.021) node[above=3pt,figlabel] {$u_{1\to2}$};
  \fill[RadialTeal] (C) circle (.021) node[below=3pt,figlabel] {$C_2$};
  \fill[CenterBlue] (Q) circle (.022) node[above=3pt,figlabel] {$q_2=1-\delta$};
  \draw[gap] (C)--(Q);
  \fill[DemandRed] (.66,1.12) circle (.022)
    node[below=3pt,figlabel] {$Z_2$};
  \node[figlabel,align=center] at (.5,.10)
    {$u_{1\to2},C_2<q_2$; choose any $Z_2$ in the displayed interval\\
     coordinate is measured from $V_2$ toward $O$};
\end{tikzpicture}%
}
\end{minipage}\hfill
\begin{minipage}[t]{.485\linewidth}
\centering
\resizebox{.96\linewidth}{!}{%
\begin{tikzpicture}[x=6cm,y=1cm,>=Latex]
  \node[panellabel,anchor=west] at (-.02,2.28)
    {(b) nonadjacent Vd role: the residual-tail point};
  \draw[flow] (0,1.12)--(1,1.12);
  \fill (0,1.12) circle (.017) node[left=2pt,figlabel] {$O$};
  \fill (1,1.12) circle (.017) node[right=2pt,figlabel] {$V_\tau$};
  \coordinate (E) at (.23,1.12);
  \coordinate (Vd) at (.39,1.12);
  \coordinate (D) at (.56,1.12);
  \draw[centertrace] (0,1.12)--(E);
  \draw[actualreach] (1,1.12)--(Vd);
  \fill[CenterBlue] (E) circle (.022) node[above=3pt,figlabel] {$T=\alpha+\delta$};
  \fill[VertexOrange] (Vd) circle (.021) node[below=3pt,figlabel] {Vd endpoint};
  \fill[DemandRed] (D) circle (.024) node[above=3pt,figlabel] {$D_\tau$};
  \draw[gap] (E)--(D);
  \node[figlabel,align=center] at (.5,.10)
    {$T<\min\{\rho_R,\rho_L\}$ and\\
     $D_\tau=\min\{\rho_R,\rho_L\}V_\tau$ lies beyond both traces};
\end{tikzpicture}%
}
\end{minipage}

\vspace{.4em}

\begin{minipage}[t]{.485\linewidth}
\centering
\resizebox{.96\linewidth}{!}{%
\begin{tikzpicture}[x=6cm,y=1cm,>=Latex]
  \node[panellabel,anchor=west] at (-.02,2.28)
    {(c) Vd1 midpoint rescue};
  \draw[flow] (0,1.12)--(1,1.12);
  \fill (0,1.12) circle (.017) node[left=2pt,figlabel] {$V_1$};
  \fill (1,1.12) circle (.017) node[right=2pt,figlabel] {$O$};
  \coordinate (C) at (.23,1.12);
  \coordinate (M) at (.50,1.12);
  \coordinate (U) at (.72,1.12);
  \draw[actualreach] (C)--(U);
  \fill[CenterBlue] (M) rectangle +(0.025,0.025);
  \node[above=3pt,figlabel] at (M) {$M_1$};
  \draw[VertexOrange,fill=white,line width=.7pt] (U) circle (.025);
  \fill[DemandRed] (U) circle (.011);
  \node[above=3pt,figlabel] at (U) {$P_{\mathrm{Vd1}}$};
  \node[figlabel,align=center] at (.5,.10)
    {the same O-side endpoint mechanism as in the T3-like case;\\
     only the local normal-form inequality changes};
\end{tikzpicture}%
}
\end{minipage}\hfill
\begin{minipage}[t]{.485\linewidth}
\centering
\resizebox{.96\linewidth}{!}{%
\begin{tikzpicture}[x=1cm,y=1cm,>=Latex]
  \node[panellabel,anchor=west] at (-.05,2.35)
    {(d) two-chart replacement when both roles are away from $T_0$};
  \node[atlasbox,align=center,figlabel,text width=2.6cm]
    (old) at (1.45,1.12)
    {original Vd1--supercritical\\pair in one vertex chart};
  \node[atlasbox,align=center,figlabel,text width=2.6cm]
    (new) at (5.15,1.12)
    {two nonsupercritical Vd0\\roles in the correct charts};
  \node[atlasbox,align=center,figlabel,text width=2.6cm]
    (rank) at (8.85,1.12)
    {recompute $N'_{\rm gap}$:\\$0$, $1$, or $2$};
  \draw[flow] (old)--node[above,figlabel] {replace} (new);
  \draw[flow] (new)--node[above,figlabel] {do not preserve input rank} (rank);
  \node[figlabel,align=center] at (8.85,.08)
    {$0$: trace length\qquad$1$: disk--gap\qquad$2$: short ray};
\end{tikzpicture}%
}
\end{minipage}
\caption{The four geometric terminals in the one-Vd placement assembly.
In the adjacent placement the radial coordinate is measured from $V_2$
toward $O$; the two local V endpoints lie before the center interval entry
$q_2=1-\delta$, leaving a literal uncovered interval on $r_2$.  The
nonadjacent placement produces the point
$D_\tau=\min\{\rho_R,\rho_L\}V_\tau$ beyond both the Vd and center traces;
the Vd1 midpoint rescue uses the same O-side endpoint construction as the
T3-like case; and the remaining placement is reduced by a two-chart
replacement followed by recomputation of the output gap rank.  Lengths are
schematic; all labelled order relations are the proof relations.}
\label{fig:one-vd-placement-terminals}
\end{figure}


\begin{corollary}[Direct one-Vd placement assembly]
\label{prop:new-one-vd-assembly}
Assume CE2, $N_+=1$, exactly one Vd1/Vd2 role, and at least one V-gap.  Then
the skeleton is not covered.
\end{corollary}

\begin{proof}
If a further role has positive adjacent support, use Method~1.  Otherwise,
if $\sigma=0$, apply Proposition~\ref{prop:readable-one-vd-adjacent} or
\ref{prop:readable-one-vd-nonadjacent}.  If $\tau=0$, Vd2 is closed by
Proposition~\ref{prop:ce2-vd2-midpoint-length}, and Vd1 by
Proposition~\ref{prop:readable-vd1-midpoint-rescuer}.  If neither index is
zero, Vd2 again uses Method~1 and Vd1 uses
Proposition~\ref{prop:readable-one-vd-replacement}.
\end{proof}

\CaseCard{NG4}{CE2, $N_+=1$, exactly one Vd1/Vd2 role.}
{A radial point, a Vd1 O-side endpoint, or a replacement output.}
{One-third radial separation, the common rescuer budget, or output rerouting.}
{Routing row R13.}

\subsection{Nonzero-gap case dispatch}
\label{subsec:nonzero-gap-dispatch}

\begingroup
\scriptsize
\renewcommand{\arraystretch}{1.08}
\setlength{\tabcolsep}{3pt}
\begin{longtable}{@{}p{.07\textwidth}p{.27\textwidth}p{.27\textwidth}p{.30\textwidth}@{}}
\caption{Reader-oriented nonzero-gap Method~3 dispatch.}
\label{tab:readable-finite-dispatch}\\
\toprule
ID&hypotheses&forced object&terminal\\
\midrule
\endfirsthead
\toprule
ID&hypotheses&forced object&terminal\\
\midrule
\endhead
NG0&one gap, $N_+=0$, all Vd0/T3-like&common radial disk and actual gap&complementary-gap theorem\\
NG1&two gaps, $N_+\in\{0,1\}$, all Vd0/T3-like&$D_2,D_4$&short $3e$--$(1-5e)$ CE2 theorem\\
NG2&one gap, $N_+=1$, all Vd0&seven-point set $K_{\rm tr}$&endpoint squeeze versus transverse return\\
NG3&one gap, $N_+=1$, exactly one T3-like&O-side endpoint $P_T$&common rescuer-tail budget\\
NG4&two gaps, $N_+=1$, exactly one T3-like&$D_2,D_4$&short CE2 theorem\\
NG5&one-Vd adjacent to supercritical $T_0$&point on $r_2$&one-third radial separation\\
NG6&one-Vd nonadjacent to supercritical $T_0$&$D_\tau$&nonadjacent radial separation\\
NG7&Vd based at $T_0$ or pair away from $T_0$&Vd1 endpoint or replacement output&rescuer, Method~1, or rerouting\\
\bottomrule
\end{longtable}
\endgroup

\begin{proposition}[All nonzero-gap finite-enclosure rows]
\label{prop:finite-enclosure-nonzero-gap-branches}
Every nonzero-gap row of Table~\ref{tab:routing} assigned to Method~3 is
impossible.
\end{proposition}

\begin{proof}
Rows R7 and R9 are Corollary~\ref{prop:new-nplus-zero-gap-closures}.  Row R10
uses Proposition~\ref{prop:new-nplus-one-all-vd0} for one gap and
Proposition~\ref{prop:readable-two-gap-common-closure} for two gaps.  Row R11
is Corollary~\ref{prop:new-one-t3-terminal}.  The finite-enclosure portion of
row R13 is Corollary~\ref{prop:new-one-vd-assembly}.
\end{proof}
